\section{Future Work}
\label{chap:conclusions:future}

The testbed implemented in this dissertation establishes a solid foundation for the development, operation and validation of network applications and its interaction with the \acs{5G} network, through the exposure of network interfaces based on the \acs{TMF} \acs{ODA}, \acs{NEF}, CAMARA and \acs{CAPIF} standards. Nevertheless, the testbed could be further improved to expand the set of capabilities available to vertical experimenters.

Future work would extend these capabilities toward closed-loop automation. In its current form, the testbed's observability layer follows a purely reactive model, notifying applications of changes in network conditions only after they occur. Integrating an \acs{NWDAF} would address this by using the metrics already collected and exposed via the \acs{NEF} as input for predictive analytics. The AnalyticsExposure and Connectivity Insights Subscriptions interfaces could then be extended with this predictive data, allowing applications to act proactively before network conditions affect or compromise their services.

The validation procedure, defined in Section~\ref{chap:proposal:validation} for applications and exposure interfaces, would also benefit from greater automation. To this end, the testbed could automatically collect network monitoring metrics, application and orchestration logs, and exposure \acs{API} responses from the services deployed across the testbed, comparing them against the expected values, such as the attributes of the ordered \acs{NEST} or the application \acsp{KPI}. This could be implemented through a \acs{CI} platform properly integrated with the \acs{OSS}, through which verticals would configure and execute the test cases for their experiments. It would also open the possibility for multi-stage testing pipelines, in which each service can be validated independently in isolation before the next one is deployed and validated alongside it.

For consistency with the rest of the testbed, such a pipeline would also rely on \acs{TMF} Open \acsp{API} and natively integrate with OpenSlice, the testbed's \acs{OSS}. Test cases could be declared and configured through the \acs{TMF} 653 \acs{API} (Service Test Management), alongside the services ordered through the \acs{TMF} 641 \acs{API} (Service Ordering)~\cite{openAPIsTMF}. The resulting verdicts would be aggregated into reports associated with the corresponding service instances and made available to vertical experimenters through OpenSlice.


focusing on automating the validation pipeline through a Continuous Integration platform.
